\section{Purpose and audience}

This part is the detailed user manual for operators who need to start, monitor, and stop the BDX slow-control prototype without modifying its software. It explains the normal Ubuntu workflow, the Phoebus operator displays, the low-voltage power supplies, the chiller, the Raspberry Pi temperature IOC, archived history, and common faults.

The main Ubuntu host runs the global, low-voltage power-supply, and chiller IOC groups. A Raspberry Pi runs the environmental temperature IOC. Phoebus is the graphical interface, while EPICS Archiver Appliance stores selected process variables for later retrieval.

\begin{warningbox}[Operational safety]
The slow-control software is not a substitute for hardwired protection, approved operating procedures, or direct verification at the equipment. Before enabling an output or changing a setpoint, confirm the intended device, channel, limits, load state, and detector condition.
\end{warningbox}

\section{System summary}

\begin{center}
\begin{tabularx}{\textwidth}{@{}l l X@{}}
\toprule
Component & Address & Function \\
\midrule
Main IOC host & \code{172.22.50.2} & Global state, LV1, LV2, and chiller IOC groups \\
Raspberry Pi & \code{172.22.50.10} & MCP9808 environmental temperature IOC \\
LV1 & \code{172.22.50.20:9221} & TTi CPX400DP low-voltage supply \\
LV2 & \code{172.22.50.21:9221} & TTi CPX400DP low-voltage supply \\
Chiller & \code{172.22.50.60:54321} & LAUDA ECO Silver RE 1225 S \\
Archiver & \code{127.0.0.1:17665--17668} & Local Archiver Appliance services \\
\bottomrule
\end{tabularx}
\end{center}

The default main-host profile is stored in \file{config/profiles/default}. The Raspberry profile is stored separately in \file{config/profiles/raspberry}. The split profile prevents the main host and Raspberry from exposing the same environment PV names.

\section{Before starting}

Run normal operator commands from a terminal opened inside the Ubuntu graphical desktop or VNC session. A plain non-graphical SSH shell cannot launch the normal Phoebus workflow.

Enter the repository:

\begin{lstlisting}[style=shell]
cd ~/SlowControl/app/bdx-slow-control
\end{lstlisting}

If an installed command is not available in the shell:

\begin{lstlisting}[style=shell]
source .venv/bin/activate
\end{lstlisting}

The local file \file{config/runtime.env} must identify the main IOC host. Create it from the template if necessary:

\begin{lstlisting}[style=shell]
cp config/runtime.env.example config/runtime.env
\end{lstlisting}

For the current deployment it must contain:

\begin{lstlisting}[style=shell]
BDX_MAIN_HOST=172.22.50.2
\end{lstlisting}

Do not use \code{127.0.0.1} for normal operation. Remote Channel Access clients cannot reach a loopback-only IOC.

\section{Starting the system}

\subsection{Raspberry environment IOC}

The main startup command checks whether \pv{BDX:ENV:TEMP:T00:VALUE} is reachable but does not start the Raspberry service automatically. Start the remote service when required:

\begin{lstlisting}[style=shell]
start-bdx-raspberry-ioc
\end{lstlisting}

Force a restart with:

\begin{lstlisting}[style=shell]
start-bdx-raspberry-ioc --restart
\end{lstlisting}

The installed command is declared in \file{pyproject.toml} and implemented by \file{src/bdx_slow_control/operator_commands.py}. It connects through SSH, runs \command{systemctl start} or \command{systemctl restart} for the remote \command{bdx-environment-ioc} service, and waits for the readiness PV. If it fails, inspect the remote service:

\begin{lstlisting}[style=shell]
ssh -t pi@172.22.50.10 \
  sudo journalctl -u bdx-environment-ioc -n 100 --no-pager
\end{lstlisting}

\subsection{Complete operator stack}

\begin{lstlisting}[style=shell]
bdx_slow_control_start overview
\end{lstlisting}

The command:

\begin{enumerate}
  \item resolves the repository and validates the main-host address, caproto client, Phoebus installation, and graphical terminal;
  \item reports the independent Archiver Appliance endpoint health in read-only mode;
  \item checks the Raspberry environment IOC and reports a warning if unavailable;
  \item opens a terminal for the main aggregated IOC unless one is already listening;
  \item waits for the IOC and a representative PV in the Phoebus terminal workflow;
  \item launches Phoebus on the requested display.
\end{enumerate}

The command never starts, stops, audits, repairs, or registers PVs in the Archiver Appliance. Use \command{bdx_archiver_start}, \command{bdx_archiver_audit}, \command{bdx_archiver_repair}, or \command{bdx_archiver_kill} explicitly when those actions are required.

The installed entry point is declared in \file{pyproject.toml} and calls \file{src/bdx_slow_control/operator_startup.py}. The Phoebus terminal workflow sources \file{scripts/start_bdx_stack.sh}; the main IOC command executed is \command{bdx-prototype-ioc}.

Startup does not enable PSU outputs, start or stop the chiller, or write setpoints. A normal launch leaves a \textbf{BDX Main IOC} terminal open when a new IOC is started and opens a separate \textbf{BDX Phoebus} terminal for the GUI workflow.

Open a different page initially with:

\begin{lstlisting}[style=shell]
bdx_slow_control_start psu
bdx_slow_control_start chiller
bdx_slow_control_start environment
bdx_slow_control_start trends
\end{lstlisting}

When the IOC is already running and only Phoebus must be reopened:

\begin{lstlisting}[style=shell]
./scripts/launch_phoebus.sh overview
\end{lstlisting}

\section{Using Phoebus}

\begin{tabularx}{\textwidth}{@{}l X@{}}
\toprule
Display & Purpose \\
\midrule
\file{overview.bob} & Overall status and navigation \\
\file{psu.bob} & Routine LV1/LV2 operation with staged requests and live plots \\
\file{psu_expert.bob} & Complete PSU PV table and direct compatibility controls \\
\file{chiller.bob} & Routine chiller operation and live temperature plots \\
\file{chiller_expert.bob} & Complete chiller PV table and expert controls \\
\file{environment.bob} & Environmental sensor status and values \\
\file{global.bob} & Global state, update timing, interlock test, reset, and all-off \\
\file{trends.bob} & Principal live and archived trends \\
\file{all_pvs.bob} & Complete configured PV catalog \\
\bottomrule
\end{tabularx}

Operator pages expose the normal workflow. Expert pages expose the complete writable PV surface and should be used only when the direct behavior is understood.

Before issuing a command, verify:

\begin{itemize}
  \item \pv{IOC_STATE} is \code{RUNNING};
  \item \pv{COMM_OK} is true;
  \item \pv{COMM_STATUS} is \code{OK} for hardware or \code{SIMULATION} for a simulator;
  \item \pv{LAST_UPDATE} continues changing;
  \item \pv{ERROR_CODE} is zero and \pv{ERROR_MESSAGE} is empty.
\end{itemize}

A visible value may be stale. If communication is unhealthy, do not assume the displayed value represents the present hardware state.

\section{Operating the low-voltage supplies}

The PSU operator page uses staged request PVs. Editing a request field does not immediately write the instrument.

For the intended device and channel:

\begin{enumerate}
  \item confirm the device label, channel, communication status, and output state;
  \item enter the target in \pv{VOLTAGE_REQUEST};
  \item enter the limit in \pv{CURRENT_LIMIT_REQUEST};
  \item check both values and press the confirmed \textbf{Apply} button;
  \item verify \pv{APPLY_STATUS} and \pv{APPLY_MESSAGE};
  \item verify \pv{VOLTAGE_SET_RBV}, \pv{VOLTAGE_RBV}, and \pv{CURRENT_LIMIT_RBV}.
\end{enumerate}

The IOC validates voltage and current together before writing either value. The default software envelope is 0--60 V, 0--20 A, and a maximum requested voltage-current product of 420 W. A profile may define stricter device or channel limits.

\begin{importantbox}[Requests and readbacks]
A request or writable setpoint indicates what was entered or accepted by the IOC. A field ending in \code{_RBV} is updated from a device query and represents the hardware readback. Always verify the readback after a command.
\end{importantbox}

Use the confirmed output control to enable or disable a channel. After the command, verify \pv{OUTPUT_RBV} and the measured voltage/current rather than relying only on the button state.

Device and global all-off commands are available in Phoebus. They act through the configured software drivers and depend on working communication. They do not replace hardware emergency-off systems.

\section{Operating the chiller}

To apply a temperature setpoint:

\begin{enumerate}
  \item verify communication and confirm that temperatures are updating;
  \item enter the target in \pv{SETPOINT_REQUEST};
  \item press the confirmed setpoint apply control;
  \item verify the apply status and applied setpoint readback;
  \item monitor controlled temperature, bath temperature, and deviation state.
\end{enumerate}

The current configured setpoint range is 5--40 degrees C. The default deviation thresholds are 0.2 degrees C for warning and 0.5 degrees C for alarm.

The chiller does not start automatically when the IOC starts. Use the confirmed \textbf{Start} or \textbf{Stop} control and verify the running-state readback and subsequent temperature evolution.

Pressure and external-temperature measurements are disabled in the current main-server profile because they are not available from the present setup. Disabled measurements are not queried and are not shown on the operator page.

\section{Environmental temperatures}

The Raspberry Pi exposes four MCP9808 temperature PVs:

\begin{lstlisting}[style=shell]
BDX:ENV:TEMP:T00:VALUE
BDX:ENV:TEMP:T01:VALUE
BDX:ENV:TEMP:T02:VALUE
BDX:ENV:TEMP:T03:VALUE
\end{lstlisting}

Read one value from the main host with:

\begin{lstlisting}[style=shell]
caproto-get BDX:ENV:TEMP:T00:VALUE
\end{lstlisting}

The detailed Raspberry installation, I2C, diagnostic, systemd, and replacement-host procedure is maintained in \file{docs/raspberry.md}.

\section{Viewing archived history}

Phoebus Data Browser combines live Channel Access data with Archiver Appliance retrieval when the PV is registered and actively sampled. Operator pages contain principal plots and \textbf{Full history} actions for archived PVs.

History exists only from the time a PV was successfully registered and sampled. Missing earlier samples cannot be reconstructed. If live values are present but history is empty, run \command{bdx_archiver_audit}, then check the retrieval service, selected time range, and clock synchronization. Use \command{bdx_archiver_repair} only when a catalog repair is actually required.

\section{Operator command implementation reference}

The installed command names are declared in \file{pyproject.toml}. Python entry points validate arguments and dispatch to the following implementation files:

\begin{longtable}{@{}p{0.27\textwidth}p{0.31\textwidth}p{0.36\textwidth}@{}}
\toprule
Command & Python entry point & Script or implementation used \\
\midrule
\endhead
\command{bdx_slow_control_start} & \file{operator_startup.py} & sources \file{scripts/start_bdx_stack.sh}; starts \command{bdx-prototype-ioc} and launches Phoebus \\
\addlinespace
\command{bdx_archiver_start} & \file{operator_startup.py} & sources \file{scripts/start_bdx_stack.sh} and calls its Archiver start/validation functions \\
\addlinespace
\command{bdx_archiver_audit} & \file{operator_startup.py} & \file{deploy/archiver-appliance/scripts/repair-archiver.sh --audit-only} \\
\addlinespace
\command{bdx_archiver_repair} & \file{operator_startup.py} & \file{deploy/archiver-appliance/scripts/repair-archiver.sh} \\
\addlinespace
\command{bdx_slow_control_kill} & \file{operator_commands.py} & \file{scripts/kill_slow_control_all.sh} \\
\addlinespace
\command{bdx_slow_control_kill_ioc} & \file{operator_commands.py} & \file{scripts/kill_slow_control_ioc.sh} \\
\addlinespace
\command{bdx_slow_control_kill_phoebus} & \file{operator_commands.py} & \file{scripts/kill_slow_control_phoebus.sh} \\
\addlinespace
\command{bdx_archiver_kill} & \file{operator_commands.py} & \file{scripts/kill_slow_control_archiver.sh} \\
\addlinespace
\command{start-bdx-raspberry-ioc} & \file{operator_commands.py} & SSH plus remote \command{systemctl start/restart bdx-environment-ioc} \\
\bottomrule
\end{longtable}

\section{Stopping the software}

Stop individual components with:

\begin{lstlisting}[style=shell]
bdx_slow_control_kill_phoebus
bdx_archiver_kill
bdx_slow_control_kill_ioc
\end{lstlisting}

Stop the normal local slow-control stack with:

\begin{lstlisting}[style=shell]
bdx_slow_control_kill
\end{lstlisting}

This stops Phoebus and the main IOC but deliberately leaves the independently managed Archiver Appliance running. Normal shutdown is graceful. Use \code{--force} only when a recognized process does not terminate normally.

\begin{warningbox}[Software shutdown does not make hardware safe]
Stopping Phoebus, Archiver Appliance, or the IOC does not switch off PSU outputs, send an automatic all-off, stop the chiller, or change setpoints. Put the equipment in the intended state before stopping the IOC and verify that state directly.
\end{warningbox}

\section{Troubleshooting reference}

\begin{longtable}{@{}p{0.30\textwidth}p{0.64\textwidth}@{}}
\toprule
Symptom & Check or action \\
\midrule
\endhead
Installed command not found & Activate the environment with \command{source .venv/bin/activate}. \\
\addlinespace
No graphical display & Run the startup command from the Ubuntu desktop or VNC session and confirm \code{DISPLAY} or \code{WAYLAND_DISPLAY} is set. \\
\addlinespace
Runtime environment missing & Create \file{config/runtime.env} from the example and set \code{BDX_MAIN_HOST=172.22.50.2}. \\
\addlinespace
Phoebus launcher missing & Set \code{BDX_PHOEBUS_HOME} or pass \code{--phoebus-home}; verify that \file{phoebus.sh} exists. \\
\addlinespace
Raspberry unavailable & Run \command{start-bdx-raspberry-ioc} and inspect the remote systemd journal if it fails. \\
\addlinespace
LV1 or LV2 unavailable & Check cabling and front-panel network configuration, then run \command{nc -vz 172.22.50.20 9221} or \command{nc -vz 172.22.50.21 9221}. \\
\addlinespace
Chiller unavailable & Run \command{nc -vz 172.22.50.60 54321} and inspect the main IOC terminal. \\
\addlinespace
PV not found & Confirm the correct Channel Access address list, the relevant IOC is running, and no duplicate IOC exposes the same PV names. \\
\addlinespace
Archiver absent during slow-control startup & Run \command{bdx_archiver_start}; \command{bdx_slow_control_start} intentionally reports Archiver health read-only and continues with live control. \\
\addlinespace
Live values but no history & Run \command{bdx_archiver_audit}; confirm retrieval availability, selected time range, and synchronized clocks; repair only when required. \\
\addlinespace
Display appears stale & Check \pv{LAST_UPDATE}, \pv{COMM_OK}, and the relevant IOC terminal. \\
\bottomrule
\end{longtable}
